\section{Постановка задачи}
\label{sec:task}

Целью лабораторной работы является изучение принципов классификации сообщений по теме, определение классов намерений и сущностей, а также формальная спецификация сценария диалога для системы NIKA.

Вариант~\mdash  15. Предметная область~\mdash  музыка (формализована в лабораторной работе №1).

Требуется:
\begin{enumerate}
	\item описать сценарий диалога не менее чем из восьми сообщений (четыре вопроса по теме и четыре ответа);
	\item обеспечить классификацию вопроса <<Что такое~\ldots?>> для формализованных сущностей и проверить ответы в диалоговом окне Ники;
	\item выделить классы и сущности сообщений сценария; каждое сообщение должно иметь хотя бы один класс;
	\item в отчёте указать цель, задание, вариант и показать тестирование сценария.
\end{enumerate}

Классификатор сообщений системы NIKA ожидает ответ в формате сервиса Wit.ai: намерение (\texttt{intent}) и сущности с ролью \texttt{rrel\_entity}. Сервис Wit.ai для новых приложений не выполняет обучение намерений~\cite{wit2851}, поэтому вместо облачного API использован локальный совместимый классификатор \texttt{witai-shim}. Он возвращает тот же JSON, читает словарь сущностей из исходников БЗ и правила намерений из файла \texttt{intents.json}.
